\section{Support-reduction proof}
\label{app:support-proof}

We expand the two finite obligations used in Theorem~1. Let $R$ be a frame with anticommuting partner $R'$ and support $W$. For each $q\in W$, write
\[
\gamma_q=(\alpha_q,\beta_q)=
(\langle R_q,R'_q\rangle,\langle S_q,R_q\rangle)\in\Ftwo^2.
\]
The first component sums to one because $R$ and $R'$ anticommute.

\begin{lemma}[Proper zero-sum subset]
If $|W|\geq3$ and $\sum_{q\in W}\alpha_q=1$, then there is a nonempty proper subset $Q\subset W$ with $|Q|\leq2$ and $\sum_{q\in Q}\gamma_q=(0,0)$.
\end{lemma}

\begin{proof}
If a coordinate has class $(0,0)$, take that singleton. Otherwise, if a class repeats, take two coordinates of that class, whose sum vanishes over $\Ftwo^2$. If neither occurs, all classes are distinct and nonzero. There are only three nonzero classes. Thus $|W|=3$ and the classes are $(0,1),(1,0),(1,1)$, whose first components sum to zero, contradicting the odd-parity assumption.
\end{proof}

Zero the letters of $R$ on $Q$. The zero sums preserve both the anticommutation bit against $R'$ and the tag syndrome, while properness keeps $R$ nonidentity. Only the chosen coordinates can change the cost. The relevant local state consists of the class $(\alpha,\beta)$, the central or noncentral multiplier, the current restoration letters in the three blocks, and the updated letters after zeroing. Complete enumeration gives 18,432 cases. In every case,
\[
\Delta\!\left(\sum_{k,q}\phi\right)\leq
\text{refunded frame-support cost};
\]
there are zero violations. The maximum net change is zero, so the exchange never increases Eq.~\eqref{eq:cost}.

Choose an unrestricted optimum with minimum total frame support. If any frame has support at least three, the lemma and local inequality produce another feasible configuration with no higher cost and strictly smaller total support, a contradiction. Therefore an optimum exists with every frame supported on at most two qubits.

For $|W|=2$, the subset lemma has four failing odd-parity patterns. This is why the argument stops at support two rather than support one; the boundary is structural, not a computational cutoff.
